\documentclass[11pt]{article}
\usepackage[margin=1in]{geometry}
\usepackage{amsmath,amssymb,amsthm,booktabs,hyperref,mathtools}
\newtheorem{theorem}{Theorem}
\newtheorem{lemma}[theorem]{Lemma}
\newtheorem{proposition}[theorem]{Proposition}
\newtheorem{conjecture}[theorem]{Conjecture}
\theoremstyle{definition}
\newtheorem{definition}[theorem]{Definition}
\newtheorem{remark}[theorem]{Remark}
\renewcommand{\Re}{\operatorname{Re}}
\title{The feature count as one inequality:\\ a local payment and a global detection}
\author{}
\date{18 October 2026}
\begin{document}
\maketitle

\begin{abstract}
The open cases of the feature-count theorem (relations, mixed slot-sets, compensation, $n\ge3$) are not different
problems.  They are one inequality, split into two ingredients.  The \emph{local} ingredient bounds how much cheaper a single
feature can be bought than its interference-free price, given its own sign-flipped coefficients.  It involves one feature's
neurons only, is finite-dimensional and does not depend on $p$.  At $n=2$ it is known sharply.  The \emph{global} ingredient
says that the interference a network carries shows up in its worst input, unless the network pays for hiding it.  All dependence
on $p$, on the frequency set, on the other features and on compensation sits in this one statement.  We state both, prove that
together they give the count for every $n,p,b$, check that every case of the earlier notes is an instance, and reduce the global
ingredient in the compensation-free class to a single statement about real trigonometric polynomials on $\mathbb Z_p^n$.  References
of the form ``Note, X'' are to \texttt{mechanism\_note.tex}.
\end{abstract}

\section{Setting}
As in the Note: $p$ an odd prime, $n\ge2$ operands, $\nu=n+1$, $b\in[2,\nu]$, $s=\nu/b$, $r=1/s$, $\rho=e^{2\pi i/p}$,
$F=\{1,\dots,\frac{p-1}2\}$, $\kappa_l(\delta)=1-\cos(2\pi l\delta/p)$, $c_l(\delta)=1-\rho^{-l\delta}$ (so $|c_l(\delta)|^2=2\kappa_l(\delta)$).
A network $\theta$ has neurons $f_i(x)[c]=W_i[c]\prod_tu_{it}(x_t)$; $P=\|\theta\|_{2,b}^b$; $g(x,\delta)=f(x)[y(x)]-f(x)[y(x)-\delta]$ with
$y(x)=\sum_tx_t$; $h(\theta)=\min_{x,\delta\ne0}g(x,\delta)$.  For $l\in F$: clean coefficient $\beta_l$ with atoms $\beta_{il}$,
$Q_l=\sum_i|\beta_{il}|^r$, interference-free (IF) price $a_l=\nu^{1-s}Q_l^{\,s}$, cheap set $L'=\{l:|\beta_l|>a_l\}$.  For $\emptyset\ne T\subseteq[n]$,
$\beta_{l,T}$ is the sign-flipped coefficient of feature $l$'s own atoms; it sits at the character $\chi_{l,T}=2l\mathbf 1_T\in\mathbb Z_p^n$
(Note, Theorem~11 and Lemma~6 of the paper: this holds for arbitrary inputs).  $\Delta(\theta)$, $\mathcal M(\theta)$, $K_l(\mu)$, the unused
budget $U(\theta)$ and the coverage gap $G(\theta)$ are as in the Note's section ``Architecture of a proof of the count''; in
particular $U+G=\nu^{1-s}(C_bP)^s\Gamma_s(p)-\min_\delta\sum_la_l\kappa_l(\delta)\ge0$.

\section{The statement}
\begin{definition}[admissible payment]
A function $\Pi=\Pi_{n,s}$ of a vector $(\beta_T)_{\emptyset\ne T\subseteq[n]}\in\mathbb C^{2^n-1}$ is an \emph{admissible payment} if for every
finite atom system $(\beta_i,\xi_{it})$ (one feature; $A=\sum_i\beta_i$, $Q=\sum_i|\beta_i|^r$, $\beta_T=\sum_i\beta_i\prod_{t\in T}\xi_{it}$)
\[(\mathrm L)\qquad \Re A-\nu^{1-s}Q^s\ \le\ \Pi\big((\beta_T)_T\big).\]
\end{definition}
By the paper's Lemma~6 every feature of every network is such an atom system, so $\Re\beta_l-a_l\le\Pi(\beta_{l,\cdot})$ for all $l$.
(L) mentions one feature only; it does not depend on $p$, the width, or the other features.

\begin{definition}[detection and the global ratio]
For a network $\theta$ with $L'(\theta)\neq\emptyset$ and $\mu\in\mathcal M(\theta)$ put
\[\mathrm{Det}_\mu(\theta)=\sum_\delta\mu_\delta\big(\mathbb E_xg(x,\delta)-\min_xg(x,\delta)\big),\qquad
\mathrm{Pay}_\mu(\theta)=\sum_{l\in L'}K_l(\mu)\,\Pi(\beta_{l,\cdot}),\]
\[R_\Pi(\theta)=\sup_{\mu\in\mathcal M(\theta)}\frac{\mathrm{Det}_\mu(\theta)+U(\theta)+G(\theta)}{\mathrm{Pay}_\mu(\theta)},\qquad
\sigma^*_\Pi(n,p,b)=\inf\{R_\Pi(\theta):L'(\theta)\ne\emptyset\}.\]
The \emph{global ingredient} is
\[(\mathrm G)\qquad R_\Pi(\theta)>1\quad\text{for every network $\theta$ with a cheap feature.}\]
\end{definition}

\begin{theorem}[the two ingredients give the count]\label{thm:main}
Fix $n,p,b$ and an admissible payment $\Pi$.  If (G) holds, then every global maximiser of the $L_{2,b}$-margin uses exactly
$|\operatorname{spt}y^\star|$ features for some maximiser $y^\star$ of $\Gamma_s(p)$.  In particular this holds if $\sigma^*_\Pi(n,p,b)>1$.
\end{theorem}
\begin{proof}
Let $\theta$ have a cheap feature.  By (G) pick $\mu$ with $\mathrm{Det}_\mu+U+G>\mathrm{Pay}_\mu$.  By (L), $\sum_l(\Re\beta_l-a_l)K_l(\mu)\le\mathrm{Pay}_\mu$
(non-cheap features contribute $\le0$, and $K_l\ge0$).  Taking point masses at the minimisers of $g(\cdot,\delta)$ as tests in the master
inequality (Note, Proposition ``master inequality'') gives the strict form of (MI), so $\gamma(\theta)<\nu^{1-s}B_p\Gamma_s(p)\le\gamma^*_b$.
Hence no global maximiser has a cheap feature, and Note, Proposition ``no cheap features $\Rightarrow$ exact count'' finishes.
\end{proof}

\begin{remark}[what each side carries]
$\mathrm{Pay}$ is linear in the features: each cheap feature contributes its own payment, weighted by how much it covers $\mu$.
$\mathrm{Det}$ is the gap between mean and worst input of the actual margin profile.  It sees every mode of $g$: the cheap features'
own interference, the other features' interference, partner and foreign modes (compensation), and all coincidences between them.
$U$ pays for norm that buys no clean amplitude (in particular for paid compensation, Note, Lemma ``paid compensation is
quadratic''), and $G$ pays for an inefficient allocation.  So (G) is a single statement about how interference, compensation and
coverage trade off.  The global ratio is measured in units of the local payment: $\sigma^*_\Pi>1$ means detection beats the best
conceivable discount.
\end{remark}

\begin{remark}[the price need not be the IF price]
Note, Proposition ``decoupling suffices'' shows that the count needs only decoupling with a \emph{common} price.  Replacing
$\nu^{1-s}Q^s$ in (L) and in $a_l$, $U$, $G$ by $\pi_{n,b}(p)Q^s$, the best single-feature price, gives a version (L$_\pi$)+(G$_\pi$) of
Theorem~\ref{thm:main} with the same proof.  Single-feature cheats are then absorbed into the price, which can only lower $\Pi$.  This is
the form to use for $n\ge3$, where the IF price may not be the true single-feature price (open problems O3, O4 of the Note).
\end{remark}

\section{Coverage: every case is an instance of (G)}
(G) quantifies over all networks, so nothing is left out.  The table records how the earlier cases sit inside it, and which parts
of (G) are proved.  ``Line'' means that all cheap features' interference sits on characters of one slot-set $T$, so that it depends on
$x$ through one linear form $s_T(x)$ only; ``plane'' (and higher) means several slot-sets.
\begin{center}
\small
\begin{tabular}{p{0.24\textwidth}p{0.30\textwidth}p{0.38\textwidth}}
\toprule
class of networks & what (G) becomes & proved\\
\midrule
interference-free & vacuous ($L'=\emptyset$) & all $n,p,b$ (Note, no-cheap-features proposition)\\
one frequency, any inputs & single cheap feature, cone geometry & $n=2$: every $p$ (Note, Thm.\ ``price of a feature''); every $n$: large $p$\\
class $\mathcal C$, line ($T=\{1,2\}$, $n=2$) & trig polynomial on $\mathbb Z_p$ with spectrum $2S$ (Prop.~\ref{prop:classC}) & $p\le59$; any $p$ if $\sum_{L'}\kappa<29.7$; near-IF; inefficient allocations\\
class $\mathcal C$, plane ($n=2$, several $T$) & trig polynomial on $\mathbb Z_p^2$ & $p=5,7$; few cheap features; (D$_3$)-dissociated cheap sets; near-IF\\
compensation (outside $\mathcal C$) & extra modes at $\chi_{l,T}$ inside $g$; paid part in $U$ & single cheap feature: $p=5,7,13$ proved, $p\le41$ numerical\\
$n\ge3$ & same, with $\Pi_n$ & $\Pi_n$ finite (crude); dissociated sets if $k^\flat(n,s)<1$\\
$p=3$ & (G) \emph{fails} & the counterexamples; see below\\
\bottomrule
\end{tabular}
\end{center}
\emph{Consistency at $p=3$.}  The certified $p=3$ networks beat the IF value, so (G) must fail there, and it does in the place the
statement predicts: their interference is hidden by compensation through the coincidence $2\equiv-1\pmod3$ (constant input
components put a mode exactly on $-\chi_{1,T}$; Note, remark after the Riesz-product theorem).  The local ingredient (L) holds at
$p=3$ (it does not depend on $p$).  So the only way the count could fail at $p\ge5$ is the same way: a global failure of detection.

\section{The compensation-free class: one statement about trigonometric polynomials}
\begin{proposition}[class $\mathcal C$]\label{prop:classC}
Let $\theta\in\mathcal C$, $a_l$ its IF prices, $e(a)=\min_\delta\sum_la_l\kappa_l(\delta)/(\Gamma_s\|a\|_{1/s})$ its efficiency, $\delta_0\in\Delta(\theta)$,
and $\mu=\frac12(\mathbf 1_{\delta_0}+\mathbf 1_{-\delta_0})$.  Then
\[\mathrm{Det}_\mu=\tfrac12\big(-\min_xF_{\delta_0}(x)-\min_xF_{-\delta_0}(x)\big),\qquad
F_\delta(x)=\sum_{l\in F}\sum_{T\ne\emptyset}|\beta_{l,T}|\,|c_l(\delta)|\cos\Big(\tfrac{2\pi}p\,2l\,s_T(x)+\psi_{l,T}+\arg c_l(\delta)\Big),\]
with phases $\psi_{l,T}$, and $G(\theta)\ge(1/e(a)-1)\sum_la_l\kappa_l(\delta_0)$, $|\beta_{l,T}|\le\nu^{s-1}a_l$.
\end{proposition}
\begin{proof}
The proof of Note, Lemma ``interference profiles are slack potentials'' gives $g(x,\delta)=\mathbb E_xg(\cdot,\delta)+F_\delta(x)$ with
$F_\delta(x)=\sum_{l,T}\Re\big(a_{l,T}c_l(\delta)\rho^{2ls_T(x)}\big)$ and $|a_{l,T}|=|\beta_{l,T}|$.  The identity for $G$ is its definition with
$\|a\|_{1/s}=\nu^{1-s}(\sum Q_l)^s$, after dropping $U\ge0$; $|\beta_{l,T}|\le\sum_i|\beta_{il}|\le Q_l^s$.
\end{proof}
So, in class $\mathcal C$, (G) follows from the following statement, in which no network appears.

\begin{conjecture}[no unicorns, class $\mathcal C$]\label{conj:unicorn}
For every allocation $a\ge0$ on $F$, every $\delta_0$ minimising $\sum_la_l\kappa_l$, every set $L'\subseteq\operatorname{spt}a$ and every family
$(\beta_{l,T},\psi_{l,T})$ with $\beta_{l,T}=0$ for $l\notin L'$, $0<\max_{l,T}|\beta_{l,T}|$ and $|\beta_{l,T}|\le\nu^{s-1}a_l$,
\[\tfrac12\big(-\min F_{\delta_0}-\min F_{-\delta_0}\big)+\big(\tfrac1{e(a)}-1\big)\sum_la_l\kappa_l(\delta_0)\ >\ \sum_{l\in L'}\kappa_l(\delta_0)\,\Pi(\beta_{l,\cdot}).\]
\end{conjecture}
At $n=2$ the admissible payments are $\Pi=c_1\max_T|\beta_T|$ with $c_1=1-3^{1-s}$ (sharp, Note, Thm.\ ``sharp pointwise payment''),
and $\Pi=c_{12}|\beta_{\{1,2\}}|$ with $c_{12}=1-(\frac23)^{s-1}$ when only $T=\{1,2\}$ is present (Note, P$_{12}$).  The conjecture is homogeneous
except for the efficiency term.  At an optimal allocation ($e=1$) it says: the worst input sees more interference than the best discount.
Two facts already in the Note show where it is delicate:
\begin{itemize}
\item \emph{Typical versus worst input.}  For $|S|=K$ frequencies the root-mean-square of $F_\delta$ is of order $\sqrt K$ times a coefficient, while
the payment is of order $K$ times a coefficient.  The conjecture therefore needs the worst of the $p^n$ inputs to see order $K$.  With
$K\approx1.4\ln p$ there are $p\approx e^{0.7K}$ inputs per line, enough for large deviations of order $K$ when the frequencies are spread out
(measured: worst input $\ge0.54K$ with adversarial phases at $p=10007$, $K=13$, against $0.18K$ needed in the line geometry).
\item \emph{Structure.}  A set with many short additive relations (a progression) admits flat polynomials (measured $0.17K$), but covers
badly ($e\approx0$).  A unicorn would be a set that is both an efficient cover and flat.
\end{itemize}

\section{Measurement protocol}
For programme supports $S=\operatorname{spt}y^\star$ at $n=2$, $b=2$ (so $e=1$), binding offsets $\delta_0$, and each geometry, compute
\[\rho_{\rm geom}(p,S,\delta_0)=\min_{(\beta_{l,T},\psi_{l,T})}\frac{\frac12(-\min F_{\delta_0}-\min F_{-\delta_0})}{\sum_{l}\kappa_l(\delta_0)\Pi(\beta_{l,\cdot})}\]
by adversarial search (smoothed minimum, random and structured restarts).  Geometries: line $T=\{1,2\}$ with $\Pi=c_{12}|\beta_{12}|$; line
$T=\{1\}$ with $\Pi=c_1|\beta_1|$ and with the numerically sharp single-slot constant; plane (all $T$) with $\Pi=c_1\max_T|\beta_T|$.  Fixing $\delta_0$
only helps the network, so each measured value is a lower bound for the corresponding restriction of $\sigma^*_\Pi$.  The conjecture at $e=1$
needs $\rho>1$.

\section{Proof plan}
\begin{enumerate}
\item \emph{Clean lemma (line).}  If the frequencies carrying interference cover $\mathbb Z_p$ with their own weights
($\min_{\delta\ne0}\sum_lv_l\kappa_l(\delta)\ge t\sum_lv_l$) and their effective number is $O(\ln p)$, then every real polynomial with those
frequencies and amplitudes has $-\min\ge\sigma(t)\sum_lv_l$ for an explicit $\sigma(t)>0$ independent of $p$: a deterministic Sudakov
minoration for short trigonometric polynomials on $\mathbb Z_p$.
\item \emph{Blind spots.}  Reduce arbitrary interference supports to the clean lemma, using the efficiency term: a sub-block that covers
badly is blind on a large set of offsets, which the rest of the support must over-cover at a cost in $e(a)$.
\item \emph{Plane to line.}  $-\min$ over $\mathbb Z_p^2$ is at least $-\min$ over any line $\{(t,mt)\}$, on which the three slot-sets become
the frequency sets $2S$, $2mS$, $2(1+m)S$.  Average or choose $m$.
\item \emph{Compensation.}  Compensating modes enter $F$ as extra coefficients at the same characters.  Paid ones are charged to $U$;
free ones need the coincidences (F1), (F2) of the Note, which are short relations of the same kind as in step 2.
\end{enumerate}
\end{document}
